\chapter{Background}
\label{chap:background}

Memory safety has become a central focus in modern programming languages due to growing concerns over the security and reliability of real-world applications. Common memory safety violations include use-after-free (accessing memory after it has been deallocated), dangling pointers (references to freed memory), and buffer overflow/underflow (reading or writing beyond allocated memory, leading to corruption of adjacent data). These violations are typically detected only at runtime and, if left unhandled, can result in critical program failures.

Memory safety describes the idea that we can reason about the program’s state locally. This allows us to show the correctness of the program component, by proving that all its memory accesses are confined to an exact region of memory assigned to it. Other regions can be ignored in this proof, as they should remain completely unaffected \cite[]{dhurjatiMemorySafetyRuntime} \cite[]{bauerPrinciplesSecurityTrust2018}.

Systems programming is a subsection of software development focused on building software and software systems that provide a service to other software. These systems often have much more rigid constraints on performance, uptime and security. Memory safe approaches are as such becoming a more interesting option when designing these systems, as they provide higher performance compared to garbage collected languages, while retaining the runtime safety these provide.

\section{Rust’s Ownership Model}

Rust replaces manual memory management with the ownership model, which is based on three fundamental rules:

1. Every value has a single owner.
\begin{lstlisting}[language=Rust, frame=single, caption={Ownership in Rust}, captionpos=b]
{
    let s = String::from("Rust"); // `s` is the owner of the String
    println!("{}", s); // `s` can be used here
} // `s` goes out of scope here
\end{lstlisting}

2. A value can have only one owner at any given time.
\begin{lstlisting}[language=Rust, frame=single, caption={Single ownership in Rust}, captionpos=b]
let s1 = String::from("Ownership");
let s2 = s1; // Ownership moves from `s1` to `s2`
println!("{}", s1); // ERROR: `s1` is no longer valid!
println!("{}", s2); // `s2` is the new owner
\end{lstlisting}

3. When the owner goes out of scope, the value is dropped \cite[Section 4.1]{RustBook}.
\begin{lstlisting}[language=Rust, frame=single, caption={Dropping ownership in Rust}, captionpos=b]
{
    let s = String::from("Temporary owner"); // `s` is created inside this block
    println!("{}", s);
} // `s` is dropped here automatically (out of scope)
println!("{}", s); // ERROR: `s` no longer exists
\end{lstlisting}

Rust also allows for borrowing or references, where multiple immutable, or a single mutable reference to a value can exist at any given time. Furthermore, all references must be valid, and invalid references are caught at compile time \cite[Section 4.2]{RustBook}. Rust’s borrow checker and ownership rules avoid the need for manual memory management and as a result ensures compile time memory safety. 

Rust extends this ownership model to threads with its built in Send and Sync traits. Traits in Rust are a collection of methods that define shared behavior, similar to interfaces in other languages. When doing concurrent programming using Rust's standard library concurrency abstractions like Mutex, the compiler checks that Send and Sync traits are valid and provides locking and unlocking guarantees at compile time \cite[Section 16]{RustBook}.

\begin{lstlisting}[language=Rust, frame=single, caption={Thread Safety using Sync and Send Traits}, captionpos=b]
let counter = Arc::new(Mutex::new(0)); // Safe shared state
for _ in 0..10 {
    let counter = Arc::clone(&counter);
    let handle = thread::spawn(move || {
        let mut num = counter.lock().unwrap(); //Compiler enforced locking
        *num += 1; // Safe modification inside Mutex
    });
}
\end{lstlisting}

Rust introduces the concept of unsafe code blocks to handle operations that bypass the compiler’s strict safety guarantees. Rust's static analysis prioritizes rejecting potentially unsafe code over allowing unsafe behavior. Unsafe blocks enable programmers to explicitly mark sections of code that the compiler would normally flag as unsafe but that the programmer deems correct. This mechanism provides controlled access to operations that Rust does not guarantee as safe, including:

\begin{enumerate}[noitemsep]
    \item Dereference a raw pointer
    \item Call and unsafe function or method
    \item Access or modify a mutable static variable
    \item Implement an unsafe trait
    \item Access fields of a union \cite[Section 20.1]{RustBook}
\end{enumerate}

For example, Rust requires the use of an unsafe block when dereferencing a raw pointer:

\begin{lstlisting}[language=Rust, frame=single, caption={Explicit unsafe declaration in Rust}, captionpos=b]
let mut num = 5;
let r1 = &raw const num;
let r2 = &raw mut num;
unsafe {
    println!("r1 is: {}", *r1);
    println!("r2 is: {}", *r2);
}
\end{lstlisting}

Rust programs often need to interact with existing functions written in languages like C or C++. For this purpose, Rust provides a Foreign Function Interface (FFI), which requires unsafe blocks because Rust cannot verify the safety of external functions, even if they are inherently safe.

\begin{table}[H]
    \centering
    \begin{lstlisting}[language=Rust, frame=single, caption={Foreign Function Interface in Rust}, captionpos=b]
unsafe extern "C" {
    safe fn abs(input: i32) -> i32;
}

fn main() {
    unsafe {
        println!("Absolute value of -3 according to C: {}", abs(-3));
    }
}
    \end{lstlisting}
\end{table}

\section{C++’s memory management model}

C/ C++ are considered memory unsafe languages compared to a memory safe language like Rust. C/ C++ allow the user to directly manage and manipulate memory, and directly manipulate raw pointers and perform pointer arithmetic. It should be noted that these operations aren’t by their nature unsafe, but can be easily implemented in an unsafe way, by design or inadvertently. 

Both C++ and Rust implement RAII (Resource Acquisition is Initialization). RAII binds the lifetime of a resource that must be acquired before use to the lifetime of the object using it. C++ does not enforce memory safety on user provided resources however, leading to potential memory management issues when working with raw pointers. The C++ standard library offers RAII wrappers to manage user provided resources in the form of smart pointers and mutexes \cite[]{RAIICppreferencecom}. Rust enforces this memory safety at compile time.
Looking at comparable C++ code we see the issues that using C++'s manual memory management causes. 

1 \& 2. A value can have multiple owners at any given time.
\begin{lstlisting}[language=C++, frame=single, caption={Single owner in C++}, captionpos=b, label={listing:unique_ptr}]
string* s1 = new string("Ownership");
string* s2 = s1; // Both `s1` and `s2` point to the same memory (no ownership enforcement)
delete s1; // Memory is freed
cout << *s1; // Undefined behavior: `s1` is deleted
cout << *s2; // Also undefined behavior: `s2` is dangling
delete s2; // Double delete! Undefined behavior
\end{lstlisting}

\begin{lstlisting}[language=C++, frame=single, caption={Single owner in C++ (using RAII wrappers)}, captionpos=b, label={listing:unique_ptr_raii}]
unique_ptr<string> s1 = make_unique<string>("Ownership");
//unique_ptr<string> s2 = s1; // ERROR: Cannot copy unique_ptr (Compiler warning)
unique_ptr<string> s2 = move(s1); // Ownership moves from `s1` to `s2`
cout << *s1; // ERROR: `s1` no longer owns the string
cout << *s2 << endl; // `s2` owns it now
\end{lstlisting}

Both the listing~\ref{listing:unique_ptr} using standard C++ and the listing~\ref{listing:unique_ptr_raii} using RAII smart pointer wrappers compile without warning using g++ and compiler flags for extensive warning logging. Both these listings experience undefined behavior during runtime and crash.

3. Manually allocated memory is not automatically freed when going out of scope.
\begin{lstlisting}[language=C++, frame=single, caption={Dropping ownership in C++}, captionpos=b, label={listing:scope_example}]
void scopeExample() {
    string* s = new string("Temporary owner"); // Manually allocated
    cout << *s << endl;
} // MEMORY LEAK: `s` is not deleted!

int main() {
    scopeExample();
    // No delete statement, memory leak occurs
}
\end{lstlisting}

\begin{lstlisting}[language=C++, frame=single, caption={Dropping ownership in C++ (using RAII wrappers)}, captionpos=b, label={listing:scope_example_raii}]
void scopeExample() {
    auto s = make_unique<string>("Temporary owner");
    cout << *s << endl;
} // Automatically deleted when going out of scope

int main() {
    scopeExample();
    return 0;
}
\end{lstlisting}

Listing~\ref{listing:scope_example} causes a memory leak that is not caught during compile or runtime. Listing~\ref{listing:scope_example_raii} fixes this by adding additional overhead with reference counting smart pointers.